%% Generated by Sphinx.
\def\sphinxdocclass{report}
\documentclass[letterpaper,10pt,russian]{sphinxmanual}
\ifdefined\pdfpxdimen
   \let\sphinxpxdimen\pdfpxdimen\else\newdimen\sphinxpxdimen
\fi \sphinxpxdimen=.75bp\relax
\ifdefined\pdfimageresolution
    \pdfimageresolution= \numexpr \dimexpr1in\relax/\sphinxpxdimen\relax
\fi
%% let collapsible pdf bookmarks panel have high depth per default
\PassOptionsToPackage{bookmarksdepth=5}{hyperref}

\PassOptionsToPackage{booktabs}{sphinx}
\PassOptionsToPackage{colorrows}{sphinx}

\PassOptionsToPackage{warn}{textcomp}
\usepackage[utf8]{inputenc}
\ifdefined\DeclareUnicodeCharacter
% support both utf8 and utf8x syntaxes
  \ifdefined\DeclareUnicodeCharacterAsOptional
    \def\sphinxDUC#1{\DeclareUnicodeCharacter{"#1}}
  \else
    \let\sphinxDUC\DeclareUnicodeCharacter
  \fi
  \sphinxDUC{00A0}{\nobreakspace}
  \sphinxDUC{2500}{\sphinxunichar{2500}}
  \sphinxDUC{2502}{\sphinxunichar{2502}}
  \sphinxDUC{2514}{\sphinxunichar{2514}}
  \sphinxDUC{251C}{\sphinxunichar{251C}}
  \sphinxDUC{2572}{\textbackslash}
\fi
\usepackage{cmap}
\usepackage[T1]{fontenc}
\usepackage{amsmath,amssymb,amstext}
\usepackage{babel}





\usepackage[Sonny]{fncychap}
\ChNameVar{\Large\normalfont\sffamily}
\ChTitleVar{\Large\normalfont\sffamily}
\usepackage{sphinx}

\fvset{fontsize=auto}
\usepackage{geometry}


% Include hyperref last.
\usepackage{hyperref}
% Fix anchor placement for figures with captions.
\usepackage{hypcap}% it must be loaded after hyperref.
% Set up styles of URL: it should be placed after hyperref.
\urlstyle{same}

\addto\captionsrussian{\renewcommand{\contentsname}{Содержание:}}

\usepackage{sphinxmessages}
\setcounter{tocdepth}{1}



\title{Project\sphinxhyphen{}weather}
\date{нояб. 01, 2025}
\release{1.0}
\author{Lamba Viktor}
\newcommand{\sphinxlogo}{\vbox{}}
\renewcommand{\releasename}{Выпуск}
\makeindex
\begin{document}

\ifdefined\shorthandoff
  \ifnum\catcode`\=\string=\active\shorthandoff{=}\fi
  \ifnum\catcode`\"=\active\shorthandoff{"}\fi
\fi

\pagestyle{empty}
\sphinxmaketitle
\pagestyle{plain}
\sphinxtableofcontents
\pagestyle{normal}
\phantomsection\label{\detokenize{index::doc}}


\sphinxAtStartPar
\sphinxstylestrong{Project\sphinxhyphen{}weather} \sphinxhyphen{} это веб\sphinxhyphen{}приложение для получения актуальной информации о погоде по названию города или географическим координатам.

\sphinxstepscope


\chapter{Модуль Project\sphinxhyphen{}weather}
\label{\detokenize{Project_weather:module-main}}\label{\detokenize{Project_weather:project-weather}}\label{\detokenize{Project_weather::doc}}\index{module@\spxentry{module}!main@\spxentry{main}}\index{main@\spxentry{main}!module@\spxentry{module}}
\sphinxAtStartPar
Точка входа для приложения Flask.
\index{index() (в модуле main)@\spxentry{index()}\spxextra{в модуле main}}

\begin{fulllineitems}
\phantomsection\label{\detokenize{Project_weather:main.index}}
\pysigstartsignatures
\pysiglinewithargsret
{\sphinxcode{\sphinxupquote{main.}}\sphinxbfcode{\sphinxupquote{index}}}
{}
{}
\pysigstopsignatures
\sphinxAtStartPar
Отображает главную страницу проекта.
\begin{quote}\begin{description}
\sphinxlineitem{Результат}
\sphinxAtStartPar
HTML\sphinxhyphen{}страница index.html

\sphinxlineitem{Тип результата}
\sphinxAtStartPar
Response

\end{description}\end{quote}

\end{fulllineitems}



\section{Модуль app.weather.controller}
\label{\detokenize{Project_weather:module-app.weather.controller}}\label{\detokenize{Project_weather:app-weather-controller}}\index{module@\spxentry{module}!app.weather.controller@\spxentry{app.weather.controller}}\index{app.weather.controller@\spxentry{app.weather.controller}!module@\spxentry{module}}
\sphinxAtStartPar
Получение погоды.
\index{find\_city\_coordinates() (в модуле app.weather.controller)@\spxentry{find\_city\_coordinates()}\spxextra{в модуле app.weather.controller}}

\begin{fulllineitems}
\phantomsection\label{\detokenize{Project_weather:app.weather.controller.find_city_coordinates}}
\pysigstartsignatures
\pysiglinewithargsret
{\sphinxcode{\sphinxupquote{app.weather.controller.}}\sphinxbfcode{\sphinxupquote{find\_city\_coordinates}}}
{\sphinxparam{\DUrole{n}{s\_city}}}
{}
\pysigstopsignatures
\sphinxAtStartPar
Поиск координат города по названию.
\begin{quote}\begin{description}
\sphinxlineitem{Параметры}
\sphinxAtStartPar
\sphinxstyleliteralstrong{\sphinxupquote{s\_city}} (\sphinxstyleliteralemphasis{\sphinxupquote{str}}) \textendash{} Название города.

\sphinxlineitem{Результат}
\sphinxAtStartPar
(широта, долгота) или None если город не найден.

\sphinxlineitem{Тип результата}
\sphinxAtStartPar
tuple

\end{description}\end{quote}

\end{fulllineitems}

\index{get\_city\_name() (в модуле app.weather.controller)@\spxentry{get\_city\_name()}\spxextra{в модуле app.weather.controller}}

\begin{fulllineitems}
\phantomsection\label{\detokenize{Project_weather:app.weather.controller.get_city_name}}
\pysigstartsignatures
\pysiglinewithargsret
{\sphinxcode{\sphinxupquote{app.weather.controller.}}\sphinxbfcode{\sphinxupquote{get\_city\_name}}}
{\sphinxparam{\DUrole{n}{lat}}\sphinxparamcomma \sphinxparam{\DUrole{n}{lon}}}
{}
\pysigstopsignatures
\sphinxAtStartPar
Получение названия города по координатам.
\begin{quote}\begin{description}
\sphinxlineitem{Параметры}\begin{itemize}
\item {} 
\sphinxAtStartPar
\sphinxstyleliteralstrong{\sphinxupquote{lat}} (\sphinxstyleliteralemphasis{\sphinxupquote{float}}) \textendash{} Широта.

\item {} 
\sphinxAtStartPar
\sphinxstyleliteralstrong{\sphinxupquote{lon}} (\sphinxstyleliteralemphasis{\sphinxupquote{float}}) \textendash{} Долгота.

\end{itemize}

\sphinxlineitem{Результат}
\sphinxAtStartPar
Название города или строка с координатами.

\sphinxlineitem{Тип результата}
\sphinxAtStartPar
str

\end{description}\end{quote}

\end{fulllineitems}

\index{get\_weather() (в модуле app.weather.controller)@\spxentry{get\_weather()}\spxextra{в модуле app.weather.controller}}

\begin{fulllineitems}
\phantomsection\label{\detokenize{Project_weather:app.weather.controller.get_weather}}
\pysigstartsignatures
\pysiglinewithargsret
{\sphinxcode{\sphinxupquote{app.weather.controller.}}\sphinxbfcode{\sphinxupquote{get\_weather}}}
{}
{}
\pysigstopsignatures
\sphinxAtStartPar
Получает данные о погоде по названию города или координатам.
\begin{quote}\begin{description}
\sphinxlineitem{Параметры}\begin{itemize}
\item {} 
\sphinxAtStartPar
\sphinxstyleliteralstrong{\sphinxupquote{city}} (\sphinxstyleliteralemphasis{\sphinxupquote{str}}\sphinxstyleliteralemphasis{\sphinxupquote{, }}\sphinxstyleliteralemphasis{\sphinxupquote{optional}}) \textendash{} Название города.

\item {} 
\sphinxAtStartPar
\sphinxstyleliteralstrong{\sphinxupquote{lat}} (\sphinxstyleliteralemphasis{\sphinxupquote{str}}\sphinxstyleliteralemphasis{\sphinxupquote{, }}\sphinxstyleliteralemphasis{\sphinxupquote{optional}}) \textendash{} Широта.

\item {} 
\sphinxAtStartPar
\sphinxstyleliteralstrong{\sphinxupquote{lon}} (\sphinxstyleliteralemphasis{\sphinxupquote{str}}\sphinxstyleliteralemphasis{\sphinxupquote{, }}\sphinxstyleliteralemphasis{\sphinxupquote{optional}}) \textendash{} Долгота.

\end{itemize}

\sphinxlineitem{Результат}
\sphinxAtStartPar
JSON с данными о погоде или ошибкой.

\sphinxlineitem{Тип результата}
\sphinxAtStartPar
Response

\sphinxlineitem{Исключение}\begin{itemize}
\item {} 
\sphinxAtStartPar
\sphinxstyleliteralstrong{\sphinxupquote{400}} \textendash{} Неверные параметры запроса

\item {} 
\sphinxAtStartPar
\sphinxstyleliteralstrong{\sphinxupquote{404}} \textendash{} Город не найден

\item {} 
\sphinxAtStartPar
\sphinxstyleliteralstrong{\sphinxupquote{500}} \textendash{} Ошибка сервера

\end{itemize}

\end{description}\end{quote}

\end{fulllineitems}

\index{get\_weather\_description() (в модуле app.weather.controller)@\spxentry{get\_weather\_description()}\spxextra{в модуле app.weather.controller}}

\begin{fulllineitems}
\phantomsection\label{\detokenize{Project_weather:app.weather.controller.get_weather_description}}
\pysigstartsignatures
\pysiglinewithargsret
{\sphinxcode{\sphinxupquote{app.weather.controller.}}\sphinxbfcode{\sphinxupquote{get\_weather\_description}}}
{\sphinxparam{\DUrole{n}{code}}}
{}
\pysigstopsignatures
\sphinxAtStartPar
Преобразует код погоды в текстовое описание.
\begin{quote}\begin{description}
\sphinxlineitem{Параметры}
\sphinxAtStartPar
\sphinxstyleliteralstrong{\sphinxupquote{code}} (\sphinxstyleliteralemphasis{\sphinxupquote{int}}) \textendash{} Код погоды от Open\sphinxhyphen{}Meteo.

\sphinxlineitem{Результат}
\sphinxAtStartPar
Описание погоды на русском языке.

\sphinxlineitem{Тип результата}
\sphinxAtStartPar
str

\end{description}\end{quote}

\end{fulllineitems}

\index{get\_weather\_from\_openmeteo() (в модуле app.weather.controller)@\spxentry{get\_weather\_from\_openmeteo()}\spxextra{в модуле app.weather.controller}}

\begin{fulllineitems}
\phantomsection\label{\detokenize{Project_weather:app.weather.controller.get_weather_from_openmeteo}}
\pysigstartsignatures
\pysiglinewithargsret
{\sphinxcode{\sphinxupquote{app.weather.controller.}}\sphinxbfcode{\sphinxupquote{get\_weather\_from\_openmeteo}}}
{\sphinxparam{\DUrole{n}{lat}}\sphinxparamcomma \sphinxparam{\DUrole{n}{lon}}}
{}
\pysigstopsignatures
\sphinxAtStartPar
Получение данных о погоде по координатам.
\begin{quote}\begin{description}
\sphinxlineitem{Параметры}\begin{itemize}
\item {} 
\sphinxAtStartPar
\sphinxstyleliteralstrong{\sphinxupquote{lat}} (\sphinxstyleliteralemphasis{\sphinxupquote{float}}) \textendash{} Широта.

\item {} 
\sphinxAtStartPar
\sphinxstyleliteralstrong{\sphinxupquote{lon}} (\sphinxstyleliteralemphasis{\sphinxupquote{float}}) \textendash{} Долгота.

\end{itemize}

\sphinxlineitem{Результат}
\sphinxAtStartPar
Данные о погоде.

\sphinxlineitem{Тип результата}
\sphinxAtStartPar
dict

\sphinxlineitem{Исключение}
\sphinxAtStartPar
\sphinxstyleliteralstrong{\sphinxupquote{Exception}} \textendash{} Ошибка получения данных о погоде.

\end{description}\end{quote}

\end{fulllineitems}

\index{get\_weather\_main() (в модуле app.weather.controller)@\spxentry{get\_weather\_main()}\spxextra{в модуле app.weather.controller}}

\begin{fulllineitems}
\phantomsection\label{\detokenize{Project_weather:app.weather.controller.get_weather_main}}
\pysigstartsignatures
\pysiglinewithargsret
{\sphinxcode{\sphinxupquote{app.weather.controller.}}\sphinxbfcode{\sphinxupquote{get\_weather\_main}}}
{\sphinxparam{\DUrole{n}{code}}}
{}
\pysigstopsignatures
\sphinxAtStartPar
Преобразует код погоды в основную категорию.
\begin{quote}\begin{description}
\sphinxlineitem{Параметры}
\sphinxAtStartPar
\sphinxstyleliteralstrong{\sphinxupquote{code}} (\sphinxstyleliteralemphasis{\sphinxupquote{int}}) \textendash{} Код погоды от Open\sphinxhyphen{}Meteo.

\sphinxlineitem{Результат}
\sphinxAtStartPar
Основная категория погоды на английском.

\sphinxlineitem{Тип результата}
\sphinxAtStartPar
str

\end{description}\end{quote}

\end{fulllineitems}

\index{weather\_page() (в модуле app.weather.controller)@\spxentry{weather\_page()}\spxextra{в модуле app.weather.controller}}

\begin{fulllineitems}
\phantomsection\label{\detokenize{Project_weather:app.weather.controller.weather_page}}
\pysigstartsignatures
\pysiglinewithargsret
{\sphinxcode{\sphinxupquote{app.weather.controller.}}\sphinxbfcode{\sphinxupquote{weather\_page}}}
{}
{}
\pysigstopsignatures
\sphinxAtStartPar
Отображает страницу с погодой.
\begin{quote}\begin{description}
\sphinxlineitem{Параметры}\begin{itemize}
\item {} 
\sphinxAtStartPar
\sphinxstyleliteralstrong{\sphinxupquote{city}} (\sphinxstyleliteralemphasis{\sphinxupquote{str}}\sphinxstyleliteralemphasis{\sphinxupquote{, }}\sphinxstyleliteralemphasis{\sphinxupquote{optional}}) \textendash{} Название города.

\item {} 
\sphinxAtStartPar
\sphinxstyleliteralstrong{\sphinxupquote{lat}} (\sphinxstyleliteralemphasis{\sphinxupquote{str}}\sphinxstyleliteralemphasis{\sphinxupquote{, }}\sphinxstyleliteralemphasis{\sphinxupquote{optional}}) \textendash{} Широта.

\item {} 
\sphinxAtStartPar
\sphinxstyleliteralstrong{\sphinxupquote{lon}} (\sphinxstyleliteralemphasis{\sphinxupquote{str}}\sphinxstyleliteralemphasis{\sphinxupquote{, }}\sphinxstyleliteralemphasis{\sphinxupquote{optional}}) \textendash{} Долгота.

\end{itemize}

\sphinxlineitem{Результат}
\sphinxAtStartPar
HTML\sphinxhyphen{}страница weather.html с переданными параметрами.

\sphinxlineitem{Тип результата}
\sphinxAtStartPar
Response

\end{description}\end{quote}

\end{fulllineitems}



\chapter{Быстрый старт}
\label{\detokenize{index:id1}}
\sphinxAtStartPar
Чтобы начать работу с приложением:
\begin{enumerate}
\sphinxsetlistlabels{\arabic}{enumi}{enumii}{}{.}%
\item {} 
\sphinxAtStartPar
\sphinxstylestrong{Установите зависимости}:

\begin{sphinxVerbatim}[commandchars=\\\{\}]
pip\PYG{+w}{ }install\PYG{+w}{ }\PYGZhy{}r\PYG{+w}{ }requirements.txt
\end{sphinxVerbatim}

\item {} 
\sphinxAtStartPar
\sphinxstylestrong{Запустите приложение}:

\begin{sphinxVerbatim}[commandchars=\\\{\}]
python\PYG{+w}{ }main.py
\end{sphinxVerbatim}

\item {} 
\sphinxAtStartPar
\sphinxstylestrong{Откройте в браузере}:

\begin{sphinxVerbatim}[commandchars=\\\{\}]
http://localhost:5000
\end{sphinxVerbatim}

\end{enumerate}


\chapter{Основные возможности}
\label{\detokenize{index:id2}}
\sphinxAtStartPar
\sphinxstylestrong{Поиск по названию города} \sphinxhyphen{} введите название города на русском или английском

\sphinxAtStartPar
\sphinxstylestrong{Поиск по координатам} \sphinxhyphen{} укажите широту и долготу

\sphinxAtStartPar
\sphinxstylestrong{Текущая погода} \sphinxhyphen{} температура, влажность, скорость ветра

\sphinxAtStartPar
\sphinxstylestrong{Прогноз на день} \sphinxhyphen{} минимальная и максимальная температура

\sphinxAtStartPar
\sphinxstylestrong{Подробное описание} \sphinxhyphen{} текстовое описание погодных условий


\chapter{API Endpoints}
\label{\detokenize{index:api-endpoints}}

\begin{savenotes}\sphinxattablestart
\sphinxthistablewithglobalstyle
\centering
\begin{tabulary}{\linewidth}[t]{TTT}
\sphinxtoprule
\sphinxstyletheadfamily 
\sphinxAtStartPar
Endpoint
&\sphinxstyletheadfamily 
\sphinxAtStartPar
Метод
&\sphinxstyletheadfamily 
\sphinxAtStartPar
Описание
\\
\sphinxmidrule
\sphinxtableatstartofbodyhook
\sphinxAtStartPar
\sphinxcode{\sphinxupquote{/}}
&
\sphinxAtStartPar
GET
&
\sphinxAtStartPar
Главная страница приложения
\\
\sphinxhline
\sphinxAtStartPar
\sphinxcode{\sphinxupquote{/weather}}
&
\sphinxAtStartPar
GET
&
\sphinxAtStartPar
JSON API для получения погоды
\\
\sphinxhline
\sphinxAtStartPar
\sphinxcode{\sphinxupquote{/weather\sphinxhyphen{}page}}
&
\sphinxAtStartPar
GET
&
\sphinxAtStartPar
Страница с интерфейсом погоды
\\
\sphinxbottomrule
\end{tabulary}
\sphinxtableafterendhook\par
\sphinxattableend\end{savenotes}


\chapter{Примеры использования API}
\label{\detokenize{index:api}}
\begin{sphinxVerbatim}[commandchars=\\\{\}]
\PYG{c+c1}{\PYGZsh{} Получить погоду по названию города}
curl\PYG{+w}{ }\PYG{l+s+s2}{\PYGZdq{}http://localhost:5000/weather?city=Москва\PYGZdq{}}

\PYG{c+c1}{\PYGZsh{} Получить погоду по координатам}
curl\PYG{+w}{ }\PYG{l+s+s2}{\PYGZdq{}http://localhost:5000/weather?lat=55.7558\PYGZam{}lon=37.6173\PYGZdq{}}
\end{sphinxVerbatim}

\begin{sphinxVerbatim}[commandchars=\\\{\}]
\PYG{p}{\PYGZob{}}
\PYG{+w}{  }\PYG{n+nt}{\PYGZdq{}status\PYGZdq{}}\PYG{p}{:}\PYG{+w}{ }\PYG{l+m+mi}{0}\PYG{p}{,}
\PYG{+w}{  }\PYG{n+nt}{\PYGZdq{}data\PYGZdq{}}\PYG{p}{:}\PYG{+w}{ }\PYG{p}{\PYGZob{}}
\PYG{+w}{    }\PYG{n+nt}{\PYGZdq{}name\PYGZdq{}}\PYG{p}{:}\PYG{+w}{ }\PYG{l+s+s2}{\PYGZdq{}Москва\PYGZdq{}}\PYG{p}{,}
\PYG{+w}{    }\PYG{n+nt}{\PYGZdq{}main\PYGZdq{}}\PYG{p}{:}\PYG{+w}{ }\PYG{p}{\PYGZob{}}
\PYG{+w}{      }\PYG{n+nt}{\PYGZdq{}temp\PYGZdq{}}\PYG{p}{:}\PYG{+w}{ }\PYG{l+m+mf}{15.5}\PYG{p}{,}
\PYG{+w}{      }\PYG{n+nt}{\PYGZdq{}temp\PYGZus{}min\PYGZdq{}}\PYG{p}{:}\PYG{+w}{ }\PYG{l+m+mf}{12.0}\PYG{p}{,}
\PYG{+w}{      }\PYG{n+nt}{\PYGZdq{}temp\PYGZus{}max\PYGZdq{}}\PYG{p}{:}\PYG{+w}{ }\PYG{l+m+mf}{18.0}
\PYG{+w}{    }\PYG{p}{\PYGZcb{},}
\PYG{+w}{    }\PYG{n+nt}{\PYGZdq{}weather\PYGZdq{}}\PYG{p}{:}\PYG{+w}{ }\PYG{p}{[\PYGZob{}}
\PYG{+w}{      }\PYG{n+nt}{\PYGZdq{}description\PYGZdq{}}\PYG{p}{:}\PYG{+w}{ }\PYG{l+s+s2}{\PYGZdq{}Ясно\PYGZdq{}}\PYG{p}{,}
\PYG{+w}{      }\PYG{n+nt}{\PYGZdq{}main\PYGZdq{}}\PYG{p}{:}\PYG{+w}{ }\PYG{l+s+s2}{\PYGZdq{}Clear\PYGZdq{}}
\PYG{+w}{    }\PYG{p}{\PYGZcb{}]}
\PYG{+w}{  }\PYG{p}{\PYGZcb{}}
\PYG{p}{\PYGZcb{}}
\end{sphinxVerbatim}


\chapter{Структура проекта}
\label{\detokenize{index:id3}}
\begin{sphinxVerbatim}[commandchars=\\\{\}]
project\PYGZhy{}weather/
├── app/
│   ├── weather/
│   │   └── controller.py    \PYGZsh{} Основная бизнес\PYGZhy{}логика
│   ├── templates/           \PYGZsh{} HTML шаблоны
│   └── static/             \PYGZsh{} CSS, JS, изображения
├── docs/                   \PYGZsh{} Документация
├── main.py                \PYGZsh{} Точка входа Flask приложения
└── requirements.txt       \PYGZsh{} Зависимости Python
\end{sphinxVerbatim}


\chapter{Используемые технологии}
\label{\detokenize{index:id4}}\begin{itemize}
\item {} 
\sphinxAtStartPar
\sphinxstylestrong{Backend}: Python, Flask

\item {} 
\sphinxAtStartPar
\sphinxstylestrong{Frontend}: HTML, CSS, JavaScript

\item {} 
\sphinxAtStartPar
\sphinxstylestrong{API}: Open\sphinxhyphen{}Meteo

\item {} 
\sphinxAtStartPar
\sphinxstylestrong{Документация}: Sphinx

\item {} 
\sphinxAtStartPar
\sphinxstylestrong{База данных}: Postgres

\end{itemize}


\chapter{Дополнительная информация}
\label{\detokenize{index:id5}}\begin{itemize}
\item {} 
\sphinxAtStartPar
{\hyperref[\detokenize{Project_weather::doc}]{\sphinxcrossref{\DUrole{doc}{Модуль Project\sphinxhyphen{}weather}}}} \sphinxhyphen{} Полная документация по проекту

\item {} 
\sphinxAtStartPar
\sphinxhref{https://github.com/ViktorLamba/project-weather}{Исходный код на GitHub}

\item {} 
\sphinxAtStartPar
\sphinxhref{https://open-meteo.com/}{Open\sphinxhyphen{}Meteo API}

\end{itemize}


\chapter{Индексы и поиск}
\label{\detokenize{index:id6}}\begin{itemize}
\item {} 
\sphinxAtStartPar
\DUrole{xref}{\DUrole{std}{\DUrole{std-ref}{genindex}}}

\item {} 
\sphinxAtStartPar
\DUrole{xref}{\DUrole{std}{\DUrole{std-ref}{modindex}}}

\item {} 
\sphinxAtStartPar
\DUrole{xref}{\DUrole{std}{\DUrole{std-ref}{search}}}

\end{itemize}


\renewcommand{\indexname}{Содержание модулей Python}
\begin{sphinxtheindex}
\let\bigletter\sphinxstyleindexlettergroup
\bigletter{a}
\item\relax\sphinxstyleindexentry{app.weather.controller}\sphinxstyleindexpageref{Project_weather:\detokenize{module-app.weather.controller}}
\indexspace
\bigletter{m}
\item\relax\sphinxstyleindexentry{main}\sphinxstyleindexpageref{Project_weather:\detokenize{module-main}}
\end{sphinxtheindex}

\renewcommand{\indexname}{Алфавитный указатель}
\printindex
\end{document}